\section{Resumo Comparativo da Parte III}

\subsection{Conclusões Principais}

\paragraph{TCL:} 
A simulação valida o Teorema do Limite Central em uma população fortemente assimétrica. A assimetria das médias amostrais diminui com $n$, aproximando-se de uma distribuição Normal. A taxa de cobertura do IC de 95\% confirma a validade de inferências baseadas na distribuição Normal, mesmo para populações não-normais.

\paragraph{RCBD:} 
O delineamento em blocos casualizados demonstra ganho de poder estatístico ao isolar a variação sistemática dos blocos. O modelo RCBD reduz o quadrado médio residual, permitindo detectar diferenças entre tratamentos com mais precisão que em um delineamento completamente casualizado.

\paragraph{ANOVA Fatorial:} 
A presença de interação significativa entre Temperatura e Pressão impede a interpretação isolada dos efeitos principais. As recomendações devem basear-se nas combinações ótimas de níveis de ambos os fatores, não em fatores individuais.

\subsection{Estrutura de Arquivos}
Os arquivos de análise foram organizados em subpastas:

\begin{itemize}
\item \texttt{analises/tcl/}: Análises do TCL
\item \texttt{analises/rcbd/}: Análises do RCBD
\item \texttt{analises/fatorial/}: Análises do Fatorial
\item \texttt{analises/resumo/}: Resumo comparativo
\end{itemize}

